% ================================================================
\renewcommand{\realmaccent}{columnStuff}  % deep violet
\chapter{Stuff}
% ================================================================

% Full-width summary table, placed immediately after the chapter title
% so it shares the chapter's opening page (see 02-body.tex for the float
% placement notes; a \begin{table*}[t] float, not \twocolumn[...]).
% REAL RULE, 2026-09-26: Stuff no longer grants Special Feats; the
% Special Feats column is gone from this table. The old Stuff Feats were
% always just items, so they now live in the catalogue below as
% Ready-Made Gadgets, bought with Stuff Points like everything else. Body, Skills,
% and Magic still grant 1 Feat at rank 3 and 2 at rank 4.
% REAL RULE, 2026-09-23: Stuff Tags column added --- a flat pool keyed to
% Stuff rank, same shape as Body Tags/Skills Known/Tags Known elsewhere in
% the book. Supersedes the older "each item grants tags equal to one less
% than its own purchased rank" rule (see Giving Your Stuff Tags below):
% that made every item's tag count depend on ITS rank, not the
% character's; this is simpler and matches how every other column works.
% REAL RULE, 2026-09-23: Max Tier column added, restating the "never buy
% above your own rank" rule from Spending Your Stuff Points as a number
% you can read straight off this table, plus the rank-4 capstone
% exception (bump one owned Tier-4 item to Tier 5) flagged with an
% asterisk --- full rule stated again in Spending Your Stuff Points below.
% REAL RULE, 2026-09-23: Stuff Points cut from 0/4/7/12 to 0/3/6/10,
% alongside cuts to Body's Stat Points and Skills Known --- smaller pools
% across the board, aimed at fewer dice to track and smaller rolls at the
% table, per playtesting. Stuff Tags unchanged.
\begin{strip}
  \begingroup\centering
  {\hdfont\Large\bfseries\color{\realmaccent} Creating Stuff\par}
  \vspace{10pt}
  \endgroup
  \begingroup
  \setlength{\tabcolsep}{6pt}\renewcommand{\arraystretch}{1.35}
  \centering\begin{tabularx}{\textwidth}{@{}>{\centering\arraybackslash}p{1.4cm} Y >{\centering\arraybackslash}p{1.5cm} >{\centering\arraybackslash}p{1.9cm} >{\centering\arraybackslash}p{1.5cm}@{}}
  \rowcolor{\realmaccent}
  {\leavevmode\hdfont\bfseries\color{white}Rank} &
  {\leavevmode\hdfont\bfseries\color{white}You Have\ldots} &
  {\leavevmode\hdfont\bfseries\color{white}Max Tier} &
  {\leavevmode\hdfont\bfseries\color{white}Stuff Points} &
  {\leavevmode\hdfont\bfseries\color{white}Stuff Tags} \\
  \rowcolor{atrBoxbg}
  \rating{1} & Nothing but the shirt on your back. & --- & \dice{0} & \dice{0} \\
  \rowcolor{white}
  \rating{2} & Basic Stuff --- armour, a weapon, a skateboard.  & \dice{2} & \dice{3} & \dice{2} \\
  \rowcolor{atrBoxbg}
  \rating{3} & Nice Stuff --- a legendary weapon, a vehicle. & \dice{3} & \dice{6} & \dice{4} \\
  \rowcolor{white}
  \rating{4} & Fancy Stuff --- a ship, a keep, incredible gear. & \dice{4}* & \dice{10} & \dice{7} \\
  \end{tabularx}
  \endgroup
  \begingroup\centering
  {\small\itshape *At rank \rating{4}, bump one owned Tier-\rating{4} item to Tier \dice{5} as a capstone --- see \emph{Spending Your Stuff Points}.\par}
  \endgroup
\end{strip}

% TODO: opening fiction for Stuff.
\begin{readaloud}
\lipsum[15][1-3]
\end{readaloud}

\chapterart{col-stuff.png}

\lettrine{S}{tuff} is what you carry from world to world. It covers your
equipment, your wealth, and the larger things you stake out and call your
own --- bases, ships, fortresses. Rank determines not just how many points
you have to spend, but the ceiling on how good any single thing you buy
is allowed to be.

% ---- Section 1: Character Creation ---------------------------------
\section{At Character Creation}

Your Stuff rank --- set when you assigned priorities at character
creation, see the table above --- grants a pool of \keyword{Stuff
Points} and a pool of \keyword{Stuff Tags}. Spend your Stuff Points on
the catalogue below and your Stuff Tags on the things that catalogue
buys.

\subsection*{Spending Your Stuff Points}
% REAL RULE, 2026-09-22: rank is not a ceiling you fill in every
% category --- it's a point budget spent across the catalogue as you
% like. A Rank 4 character doesn't automatically own a rank-4 weapon AND
% vehicle AND armour AND fortress; they own whatever combination their
% points cover.
% REAL RULE, 2026-09-23: closes a gap in the point-buy math --- nothing
% previously stopped a low-rank character from saving every point for
% one purchase above their station (Stuff 2, 4 points, one rank-4 item).
% The "You Have..." column in the table above already implied a ceiling
% ("Basic Stuff" at rank 2); this makes it an actual rule instead of an
% unenforced suggestion.
Every catalogue entry lists a rank from \rating{1} to \rating{4} ---
that number is also its \keyword{point cost}. Spend your Stuff Points
across any combination of entries, from any category, in any
combination, with one hard limit: \emph{you may never buy anything
ranked higher than your own Stuff rank.} A Stuff-\rating{2} character
can own one Basic-tier thing and a Tier-\rating{1} Gadget, or three
Tier-\rating{1} Gadgets --- never a single Fancy-tier item bought on
credit.

\begin{example}
Kira is Stuff \rating{3}: \dice{6} points, and nothing she buys can be
rank \rating{4}. She could buy a \emph{Weapon} (\dice{3}) and a
\emph{Vehicle} (\dice{3}), spending every point on two Nice-tier things
--- or trade both for three Basic-tier things instead: a second
\emph{Weapon} (\dice{2}), \emph{Armour} (\dice{2}), and one more
(\dice{2}). Same rank, same ceiling, different loadout.
\end{example}

\subsection*{The Rank-4 Capstone: Tier 5}
% REAL RULE, 2026-09-23: the one deliberate exception to "never buy above
% your own rank" and to the game's usual 1-4 rank ceiling entirely --- a
% single Tier-5 item exists nowhere else in the book. Free (no additional
% Stuff Points), one item only, and it must already be a Tier-4 item you
% own rather than a fresh purchase.
At Stuff \rating{4}, pick one Tier-\rating{4} item you already own and
mark it \keyword{Tier \dice{5}} instead --- a capstone upgrade, free of
any additional point cost. This is the one item in your possession (and,
short of a GM exception, the one item in the entire game) that reaches
above rank \rating{4}: the single sword that ends duels, the one ship
every pirate has heard of. Everything else you own is still capped at
your own rank.

\begin{example}
Kira reaches Stuff \rating{4} and already owns a Tier-\rating{4} Weapon.
She marks it Tier \dice{5} instead of buying something new --- it costs
her nothing extra, and it is now the sharpest blade she will ever carry.
\end{example}

\subsection*{Giving Your Stuff Tags}
% REAL RULE, 2026-09-23: reworked from a per-item formula ("each item
% grants tags equal to one less than its own rank") to a flat pool keyed
% to the character's Stuff rank, matching the shape of Body Tags/Skills
% Known/Tags Known elsewhere in the book --- see the note on the
% chapter-opening table above for why.
Your Stuff rank also grants a pool of \keyword{Stuff Tags} --- \dice{0}
at rank \rating{1}, \dice{2} at rank \rating{2}, \dice{4} at rank
\rating{3}, \dice{7} at rank \rating{4} --- to distribute across
whatever you own, in any combination. Pick from the example tags listed
under that kind of item, or invent your own with your GM's approval.
These work exactly like a Body Tag or a Skill --- when one clearly
applies, add one die to your pool.

\begin{example}
Kira is Stuff \rating{3}, so she has \dice{4} Stuff Tags to give out. She
puts \atrtag{item}{Fast} and \atrtag{item}{All-Terrain} on her Vehicle,
then \atrtag{item}{Concealable} and \atrtag{item}{Piercing} on her
Weapon --- two things she owns, each with a distinct personality.
\end{example}

% ---- Section 2: Rules -----------------------------------------------
\section{Stuff in Play}

\subsection*{Using Your Stuff's Tier as a Base Stat}
% REAL RULE, 2026-09-23: applies "Stats from Objects" (How to Play) to
% the Stuff catalogue specifically. Not a new rule --- the underlying
% mechanic ("take the higher of the two") is already established there;
% this section exists so a player reading Stuff in isolation doesn't have
% to go hunting for it.
% REAL FIX, 2026-09-23: the "Worked Examples" table (Item/Situation/Base
% columns) was confusing in practice --- a Tier-5 row read as ambiguous
% about whether it meant the item was Tier 4 or Tier 5, and cramming a
% whole scenario into one cell made every row need its own re-parsing.
% Replaced with a single explicit formula plus one worked example instead
% of nine compressed ones.
\begin{center}
{\Large\itshape Base Stat = Max( Item Tier, Character Stat )}
\end{center}

When you act \emph{through} something you own --- swinging a weapon,
driving a vehicle, working from inside a fortress --- compare that
item's own \keyword{tier} to whichever base Stat the action would
normally use, and roll with whichever number is bigger. This is exactly
the rule described in \emph{Stats from Objects} in \emph{How to Play},
applied to your own gear: an item can only ever help, never hurt --- it
can lend you its tier when that number beats your Stat, but it never
forces you down to a worse number than the Stat you already have.

\begin{example}
Rooster's Body is \dice{2}, but he's swinging a Tier-\rating{4} sword.
The sword's tier beats his own Body, so he rolls with a base of \dice{4}
instead of \dice{2} --- the same trade a rank-\rating{3} car makes for a
rank-\rating{2} driver in \emph{How to Play}.
\end{example}

\subsection*{Re-Kitting on Arrival}

The catalogue is deliberately \keyword{generic}. What you own is a
\emph{kind} of thing --- a vehicle, a weapon, a base --- and you
decide what that thing actually \emph{is} each time you enter a new realm.
A \emph{Weapon} might be a flaming sword under a fantasy sun and a
plasma rifle in the sprawl; it is the same line on your sheet either way,
and its tags travel with it unchanged. Each genre chapter lists suggested
translations, so you always have a world-appropriate version ready to
hand.

\subsection*{Losing Your Stuff}
% REAL RULE, 2026-09-23: what happens when Stuff is lost, broken, stolen,
% or otherwise taken off the sheet mid-world. Deliberately simple and
% player-friendly, matching the game's overall "fail fast, keep playing"
% posture --- losing gear should sting in the moment (the GM is free to
% make narrative hay of it) but never leave a character crippled for the
% rest of a world, and never requires bookkeeping across worlds.
Gear gets smashed, stolen, or left behind --- that is what makes losing
it interesting in the moment. It is never interesting as bookkeeping.
Whatever a scene costs you, everything you own comes back exactly as
built, at full Tier, the next time you cross into a new realm: your
Stuff Points and Stuff Tags were always a record of what you
\emph{chose}, at character creation, not a fixed pile of possessions
that depletes as it gets used up.

\begin{example}
Rooster's Vehicle goes off a cliff mid-chase in the sprawl. For the rest
of that world he's on foot --- a real, lasting problem the GM can lean
on. The moment the party crosses into the next realm, that same Vehicle
reappears, re-skinned to fit, exactly as tough and fast as he bought it.
\end{example}

% ---- Section 3: Options ---------------------------------------------
\section{Lifestyle \& Cash}
% REAL RULE, 2026-09-23: Lifestyle moves from a bought gearlist entry to
% an automatic grant keyed directly to Stuff rank --- it was already
% describing the same quality tiers as the "You Have..." column in the
% chapter-opening table, so charging points for it was double-counting.
% Everything else in the chapter still costs points; this is the one
% exception, and the text below says so explicitly. Placed BEFORE "What
% You Can Buy" (moved 2026-09-23) so the one free, no-roll-required grant
% reads before the point-buy catalogue starts, not folded into it.
\artinline{stuff-lifestyle.png}
Cash in your pocket, clothes on your back, ordinary tools within reach
--- every traveller has \emph{some} baseline standard of living, and
unlike everything else in this chapter, you never spend a point on it.
Your Lifestyle simply matches your Stuff rank.

\begin{statblock}[Lifestyle \& Cash, by Rank]
\begin{tabularx}{\linewidth}{@{}p{0.16\linewidth} Y@{}}
\rating{1} & Scraping By --- basic clothes, pocket change, nothing to spare.\\[3pt]
\rating{2} & The Basics --- some cash, appropriate clothes, and basic tools. You get by.\\[3pt]
\rating{3} & Well Off --- well dressed, generally nice things, spending money, good tools.\\[3pt]
\rating{4} & Rich --- the finest version of every tool you need, fine clothes, and a wallet flush with cash.\\
\end{tabularx}
\end{statblock}

% REAL, 2026-09-23: this divider marks the start of the point-buy
% catalogue, so it gets chapter-title-scale treatment (\Huge, same rule
% motif as \chapter itself, see \titleformat{\chapter} in atrpg.sty)
% instead of an ordinary \section --- deliberately bigger than every
% other heading in the chapter. Title and banner are ONE figure* float
% (not \strip, and not \artbanner immediately following a separate
% \strip) --- \strip assumes the two columns are height-balanced at the
% point it fires, which broke here: Lifestyle & Cash's inline image
% (\artinline, not a float) leaves the two columns visibly uneven just
% above this point, so a bare \strip landed at the shorter column's
% height and overlapped the still-rendering taller one. A single figure*
% carrying both the title and the image is placed by LaTeX's normal float
% algorithm instead, the same mechanism that already placed a bare
% \artbanner correctly here before Lifestyle moved up --- immune to
% column imbalance by construction.
% REAL BUG FIX, 2026-09-23: a bare figure*[tb] is still only a QUEUED
% float --- LaTeX is free to defer it past later text if it doesn't fit
% where it's first tried, and Vehicles (which follows immediately in the
% source) was winning that race, rendering on a page BEFORE this banner.
% \FloatBarrier (placeins, see atrpg.sty) forces every float queued
% before it --- this one --- to be placed before any later material is
% typeset, guaranteeing reading order at the cost of possible whitespace
% (accepted trade-off; revisit if the gap looks bad once art is final).
\begin{figure*}[tb]
\centering
{\hdfont\bfseries\color{\realmaccent}\Huge\MakeUppercase{What You Can Buy}\par}
\vspace{6pt}
{\color{atrRule}\titlerule[1pt]}
\vspace{10pt}

\includegraphics[width=\textwidth,height=0.38\textheight,keepaspectratio]{vehicle-steed.png}
\end{figure*}
\FloatBarrier

Every entry below is tagged with its point cost, shown as pips
(\rating{1}--\rating{4}). Spend your Stuff Points from the table at the
start of this chapter here, in whatever combination you like, and never
above your own rank.

\section{Vehicles}
\artinline{stuff-vehicles.png}
A vehicle carries you, your party, and your gear across the realm --- and,
with a \keyword{Persistence Point}, from one realm into the next without
changing shape. A Vehicle is a Vehicle at any rank --- what changes is
how fast, tough, and capable it is.

\gearentry{Vehicle}{2}{The humblest way to travel faster than your feet.
It gets you there, carries a little cargo, and asks little in return. No
armour to speak of and no room for the whole party, but it is yours, and
it is quick.\par\emph{Examples:} Horse, Motorcycle, Hoverboard.}

\gearentry{Vehicle}{3}{A proper conveyance with room for a few
companions and their kit. It shrugs off rough roads, hauls real cargo, and
can take a knock without falling apart. This is the workaday vehicle of the
realm --- reliable, unremarkable, and exactly enough.\par\emph{Examples:} Covered Wagon, Family Car, Cargo Skiff.}

\gearentry{Vehicle}{4}{Fast, tough, and built to turn heads or win
fights. A rank-\rating{4} Vehicle outruns, out-hauls, or out-fights
anything short of a named threat. It carries the whole party in comfort
or in armour, and it changes what encounters are even possible.\par\emph{Examples:} War-Chariot, Race Car, Gunship.}

% REAL, 2026-09-26: example tables and per-tag descriptions replaced with
% inline "Examples:" lines and a flat Example Tags list, in every category.
\subsection*{Example Tags for Vehicles}
{\raggedright\atrtag{item}{Fast}, \atrtag{item}{Agile}, \atrtag{item}{Armoured},
\atrtag{item}{Rugged}, \atrtag{item}{All-Terrain}, \atrtag{item}{Spacious},
\atrtag{item}{Amphibious},
\atrtag{item}{Airborne}, \atrtag{item}{Silent}, \atrtag{item}{Armed},
\atrtag{item}{Luxurious}\par}

\section{Weapons}
\artinline{stuff-weapons.png}
A weapon is whatever you reach for when a fight turns physical --- blade,
gun, or something stranger. A Weapon is a Weapon at any rank --- what
changes is how much of an edge it gives you.

\gearentry{Weapon}{2}{A dependable killer, nothing more and nothing
less --- a sword, a pistol. It does exactly what a weapon is supposed to
do.\par\emph{Examples:} Sword, Pistol, Sidearm.}

\gearentry{Weapon}{3}{Something with a trick to it --- a flaming blade, a
grenade launcher, a weapon that does more than simply cut or shoot.\par\emph{Examples:} Enchanted Blade, Shotgun, Plasma Cutter.}

\gearentry{Weapon}{4}{A weapon out of story --- world-shaping, named, and
feared. Everyone in the room knows what it means when you draw it.\par\emph{Examples:} Legendary Blade, Signature Sidearm, Superweapon.}

\subsection*{Example Tags for Weapons}
{\raggedright\atrtag{item}{Ranged}, \atrtag{item}{Reach}, \atrtag{item}{Rapid-Fire},
\atrtag{item}{Returning}, \atrtag{item}{Concealable}, \atrtag{item}{Silent},
\atrtag{item}{Piercing}, \atrtag{item}{Explosive}, \atrtag{item}{Elemental},
\atrtag{item}{Disarming}\par}

\section{Armour}
\artinline{stuff-armour.png}
Whatever stands between you and the worst a fight can do. Higher rank
doesn't just mean thicker plating --- it can turn a killing blow into a
bad day.

\gearentry{Light Armour}{2}{Padding and leathers --- turns aside the
first, glancing blow.\par\emph{Examples:} Leather Jerkin, Reinforced Jacket, Impact-Weave Bodysuit.}

\gearentry{Heavy Armour}{3}{Plate, a ballistic vest, a hardsuit ---
real, worn protection.\par\emph{Examples:} Plate Harness, Riot Gear, Exo-Frame.}

\gearentry{Master Armour}{4}{Armour of legend or high craft, protection
few can match.\par\emph{Examples:} Dragon-Scale Mail, Combat Suit, Void-Rated Armour.}

\subsection*{Example Tags for Armour}
{\raggedright\atrtag{item}{Impact-Resistant}, \atrtag{item}{Fireproof},
\atrtag{item}{Sealed}, \atrtag{item}{Insulated}, \atrtag{item}{Flexible},
\atrtag{item}{Lightweight}, \atrtag{item}{Discreet}, \atrtag{item}{Imposing},
\atrtag{item}{Self-Repairing}, \atrtag{item}{Anchored}\par}

\section{Locations}
\artinline{stuff-locations.png}
A place you can call your own --- somewhere to retreat to, resupply, and
plan your next move from. Like everything else in this chapter, a
Location re-skins to fit wherever you currently stand; see
\emph{Re-Kitting on Arrival}.

\gearentry{Base}{3}{A safe house, a wizard's tower, a hideout you can
return to.\par\emph{Examples:} Hidden Tower, Safehouse, Satellite Relay.}

\gearentry{Fortress}{4}{A stronghold --- walls, staff, and standing over
the land around it.\par\emph{Examples:} Castle, Corporate Compound, Orbital Station.}

\subsection*{Example Tags for Locations}
{\raggedright\atrtag{item}{Fortified}, \atrtag{item}{Hidden}, \atrtag{item}{Watchful},
\atrtag{item}{Connected}, \atrtag{item}{Stocked}, \atrtag{item}{Workshop},
\atrtag{item}{Welcoming}, \atrtag{item}{Staffed}, \atrtag{item}{Escape Route},
\atrtag{item}{Commanding}\par}

\section{Gadgets}
% REAL RULE, 2026-09-23: the deliberate catch-all category. Every other
% catalogue entry has a fixed name and blurb decided in this book;
% Gadgets are the opposite --- the player and GM design the thing itself
% at the table, and only its rank (hence its point cost and power
% ceiling) is fixed by these rules. Keeps the catalogue from needing an
% entry for every one-off tool a player might dream up.
% REAL RULE, 2026-09-26: the former Stuff Feats are merged in here as
% ordinary Gadgets --- same one-thing shape, same Tier/point cost, same
% Stuff Tag pool.
\artinline{stuff-gadgets.png}
Not everything worth owning fits neatly into a Vehicle, a Weapon, a suit
of Armour, or a Location. A \keyword{Gadget} is whatever's left ---
purpose-built gear that does one specific, useful thing. Buy one of the
Gadgets below, or design your own together with your GM when
you buy it, the same way you'd invent a Body Tag or a Skill that isn't
on the printed list.

\subsection*{Designing a Gadget with Your GM}
Pick the rank you can afford --- never above your own Stuff rank, same
as everything else in this chapter --- then agree with your GM on the
one thing it does. The rank sets the ceiling on how impressive that one
thing is allowed to be: a Tier-\rating{1} Gadget solves a single mundane
problem outright, while a Tier-\rating{4} Gadget can do something no
ordinary tool has any business doing. If you and your GM can't agree on
what a rank should buy, undershoot --- a Gadget that's a little less
flashy than its rank allows is far easier to live with at the table than
one that quietly turns out to be a spell wearing a disguise.

A Tier-\rating{1} Gadget never supplies a better base Stat --- your own
Stats start at \dice{1} --- so its value is simply that you have it,
plus whatever Stuff Tags you hang on it.

\begin{example}
Kira wants a Tier-\rating{1} Gadget. She and her GM land on a
\emph{grappling hook}: fires a hook and cable a few dozen feet, hauls her
up a wall or yanks a target off their feet. One good trick, solving
exactly one kind of problem --- getting somewhere a hook and cable can
reach.
\end{example}

% REAL RULE, 2026-09-26: the former Stuff Feats (Bottomless Bag, Cloaking
% Suit, Companion, etc.) are folded into this one alphabetical list. Their
% Tiers are a first pass (formerly free, rank-3/4-only picks) and want
% playtesting.
% REAL RULE, 2026-09-26: Tier-1 exists for Gadgets only (Bottomless Bag,
% Grappling Hook, Letters of Marque, Lockpick Rig, Translator) --- mainly
% so a Stuff-2 character's 3 points don't strand one. Vehicles, Weapons,
% Armour, and Locations deliberately stay at 2+ so they feel special.
\artinline{feat-stuff-pocket-refuge.png}

\gearentry{Bottomless Bag}{1}{Reach in and pull out any ordinary item
that could fit in the bag: rope, a lantern, a crowbar, a change of
clothes. Nothing valuable, nothing magical, nothing you couldn't buy at
a market.}

\gearentry{Cloaking Suit}{3}{While you wear it and move slowly, you are
invisible, for up to a minute at a time. After that it needs an hour to
recharge. Running or fighting makes you flicker back into view.}

\gearentry{Companion}{3}{A loyal companion with a mind of its own: a
hound, a squire, a droid. It follows you and helps where it can, and it
can be hurt, lost, or captured.}

\gearentry{Crew}{3}{A handful of people who follow your orders: sailors,
thugs, technicians. They are not heroes, and they break when things go
badly.}

\gearentry{Environment Suit}{3}{While you wear it, you can breathe
anywhere: underwater, in smoke, in poison gas, in vacuum. Heat, cold,
pressure, radiation, and disease cannot reach you. It is bulky and
obvious, and suits can be punctured.}

\gearentry{Flight Rig}{3}{A jetpack, grav-boots, or a winged harness.
While you wear it, you fly. It is loud and obvious, and it handles
poorly: tight turns and sudden stops are rolls.}

\gearentry{Grappling Hook}{1}{Fires a hook and cable a few dozen feet;
hauls you up, or hauls something else down. One good trick: getting
somewhere a hook and cable can reach.\par\emph{Examples:} Enchanted
Climbing Rope, Grapple Gun, Tractor-Line Launcher.}

\gearentry{Informant Network}{3}{In any city, after a day, you have eyes
and ears reporting to you: who has arrived, who is asking about you,
what is for sale. Outside cities, it goes quiet.}

\gearentry{Letters of Marque}{1}{Papers granting you legal authority.
Officials treat you as having every right to be where you are. Outlaws
and rival powers don't care.}

\gearentry{Lockpick Rig}{1}{A compact set of picks and probes, or the
digital equivalent, that opens what it isn't supposed to.\par\emph{Examples:}
Thieves' Tools, Lock-Bypass Kit, Cipher Wand.}

\gearentry{Master Key}{3}{Opens any lock, mechanical or electronic,
given time: a moment for a door, an hour for a vault. It always leaves
obvious signs of tampering, and magical wards stop it cold.}

\gearentry{Nanite Patch Kit}{4}{Closes a wound, mends a break, or purges
a poison almost as fast as you can apply it.\par\emph{Examples:}
Phoenix-Down Poultice, Trauma Injector, Medical Nanite Swarm.}

\gearentry{Patron}{3}{A powerful backer who takes your call: a duke, a
corporation, a temple. They help when it suits them, and they will ask
favours in return.}

\gearentry{Pocket Refuge}{3}{A door (chalked on a wall, a folding tent,
a hatch) that opens into a small, safe room only you and your guests can
enter. Whoever holds the door holds the room.}

\gearentry{Scanner}{3}{Point it at something and see through walls,
crates, and skin on its display. It takes a moment, anyone watching can
see what you are doing, and enough metal or stone blocks it, at the GM's
discretion.}

\gearentry{Signal Jammer}{3}{Blinds a device, a ward, or a whole room to
outside detection for as long as you keep it running.\par\emph{Examples:}
Rune of Silence, RF Jammer, Scrambler Emitter.}

\gearentry{Translator}{1}{While it is on you, you understand and are
understood in any spoken language, human or not. There is a lag, and
subtlety gets lost.}

\todo{Gadget ranks for the former Stuff Feats are a first pass --- they
were free Stuff Feats until 2026-09-26. Playtest the costs.}

\subsection*{Example Tags for Gadgets}
{\raggedright\atrtag{item}{Adaptive}, \atrtag{item}{Antitoxin},
\atrtag{item}{Bone-Mending}, \atrtag{item}{Compact},
\atrtag{item}{Hard-to-Trace}, \atrtag{item}{Long-Range},
\atrtag{item}{Portable}, \atrtag{item}{Precise}, \atrtag{item}{Quick-Firing},
\atrtag{item}{Quiet}, \atrtag{item}{Rapid}, \atrtag{item}{Retractable},
\atrtag{item}{Reusable}, \atrtag{item}{Selective}, \atrtag{item}{Strong Line},
\atrtag{item}{Wide-Area}\par}
